% =========================================================
% Proof Stub: Bounded Does Not Imply Convergent
% Source: volume-iii/analysis/sequences/notes/sequences/notes-convergence.tex
% Mode: standard
% =========================================================

\clearpage
\phantomsection
\hypertarget{proof-RA-ABB-P-bounded-not-conv}{}

\subsubsection[Bounded Does Not Imply Convergent]{Proof --- Bounded Does Not Imply Convergent}

\bigskip

\noindent
\textbf{Source.}
See \texttt{notes-convergence.tex}.

\vspace{0.75em}

\noindent
\textbf{Goal.}
The converse of ``convergent $\Rightarrow$ bounded'' is false: there exist bounded sequences that do not converge.

\vspace{0.75em}

\noindent
\textbf{Logical form.}
\[
  \exists (x_n) \text{ bounded},\; (x_n) \text{ does not converge}.
\]

\vspace{0.75em}

\noindent
\textbf{Proof strategy.}
Exhibit a counterexample: $(x_n) = ((-1)^n)$ is bounded (by 1) but oscillates between $-1$ and $+1$, so has no limit.

\vspace{0.75em}

\noindent
\textbf{Proof.}
\begin{proof}
\Claim The converse of ``convergent $\Rightarrow$ bounded'' is false: there exist bounded sequences that do not converge.

\Given 

\Goal 

\Strategy Exhibit a counterexample: $(x_n) = ((-1)^n)$ is bounded (by 1) but oscillates between $-1$ and $+1$, so has no limit.

\medskip

% --- define x_n = (-1)^n ---
% --- bounded: |x_n| = 1 ≤ 1 for all n ---
% --- diverges: subsequence (x_{2n}) → 1 and (x_{2n+1}) → -1; limits differ ---
% --- cite: if convergent then all subsequences share the same limit ---

\end{proof}

\begin{remark*}[Proof shape]
Counterexample. The sequence $(-1)^n$ is bounded but diverges (two subsequences have different limits).
\end{remark*}

\begin{remark*}[Dependencies]
\begin{itemize}
  \item Uniqueness of limits (Theorem~\ref{thm:unique-limits}).
  \item Subsequences inherit limits (if convergent).
\end{itemize}
\end{remark*}
